\documentclass[12pt]{article}
\usepackage[utf8]{inputenc}
\usepackage[T1]{fontenc}
\usepackage{hyperref}
\title{Zimakani Language Resources for K-12 Curriculum Development}
\author{Compiled by Claude Sonnet 5.5 (claude-sonnet-5-5) for Andr\'as Kornai}
\date{}
\begin{document}
\maketitle

\section{Introduction}
Zimakani (ISO 639-3: zik; Glottocode: zima1244) is a Papuan language of the Western Province of Papua New Guinea, spoken near the southern end of Lake Murray. Published estimates are about 1,500 speakers, a figure from 1990, and Joshua Project reports a larger figure for the people group. Dialects named are Bagwa and Dea. Tok Pisin serves wider communication. Searches found very little material, and no school materials in Zimakani were found.

\section{Primary Educational Resources}
No primers, textbooks or school programmes in Zimakani were found in search results.

\section{Linguistic and Reference Materials}
\begin{itemize}
\item Wikipedia, Zimakani language (overview, classification, mention of early gospel texts): \url{https://en.wikipedia.org/wiki/Zimakani_language}
\item Glottolog entry for Zimakani: \url{https://glottolog.org/resource/languoid/id/zima1244}
\item Trans New Guinea database page, Zimakani (South Boazi): \url{https://transnewguinea.org/language/zimakani-south-boazi}
\item Numeral-system data sheet for Zimakani: \url{https://lingweb.eva.mpg.de/channumerals/Zimakani.htm}
\end{itemize}

\section{Religious and Community Resources}
Wikipedia reports that the Unevangelized Fields Mission produced gospel texts in the language in 1956 and 1966. I did not find these texts online.
\begin{itemize}
\item Joshua Project, Bagwa Zimakani people group: \url{https://www.joshuaproject.net/people_groups/10563}
\item Joshua Project, language page for zik: \url{https://www.joshuaproject.net/index.php/languages/zik}
\end{itemize}

\section{Digital and Audio Resources}
\begin{itemize}
\item Global Recordings Network, Zimakani language page: \url{https://glorecnet-website.k.globalrecordings.net/en/language/18716}
\end{itemize}
Search results also indicated a Paradisec collection of Zimakani materials, but I did not see its URL, so it is not cited.

\section{Curriculum Development Context}
Little is available. Any curriculum would need new work: orthography development, community consultation, and documentation. It could draw on the few materials above and on Tok Pisin and English as bridge languages.

\end{document}